\begin{abstract}
Admission control decides whether an arriving request enters an LLM serving
system, and published policies differ in the quantity they read: queue depth,
KV-cache pressure, deadline slack, tenant credit. A systematic review of
fifty-five works found that none of sixteen admission-primary papers reports the
range that quantity took during its own evaluation. The consequence is specific.
A policy whose threshold was never approached and a policy that approached it
and declined to reject produce identical admission counts and identical
downstream latencies, so no reported number separates a policy that was
exercised from one that was merely present.

We measure that range directly. We define a signal-liveness probe, run two
published policies on unmodified vLLM across arrival regimes spanning the
engine's rated capacity, and report the observed signal range alongside the
admission decisions it produced. We find that liveness is a property of the
deployment configuration rather than of the policy: the same policy reading the
same quantity is exercised under one model-to-memory-to-concurrency ratio and
unreachable under another, with nothing in either paper's reported numbers
distinguishing the two. We release the harness and propose a one-column addition
to results reporting that makes the distinction checkable.
\end{abstract}
